\chapter{مبحث الناسخ والمنسوخ ومكانة السنة}

\section{مقدمة في تعريف النسخ وحكمته}
الحمد لله والصلاة والسلام على رسول الله، أما بعد؛ فقد شرع الإمام الشافعي رحمه الله تعالى في بيان مبحث عظيم من مباحث الأصول، وهو "الناسخ والمنسوخ".

\textbf{النسخ في اللغة:} يأتي بمعنى الإزالة والنقل، فيقال: نسخت الشمس الظل (أي أزالته)، ونسخت الكتاب (أي نقلته). 
\textbf{وفي اصطلاح الأصوليين:} هو رفع حكم متقدم بدليل متراخٍ عنه. والأصل فيه قول الله جل وعلا: ﴿مَا نَنَسَخْ مِنْ آيَةٍ أَوْ نُنسِهَا نَأْتِ بِخَيْرٍ مِّنْهَا أَوْ مِثْلِهَا﴾.

وقد مهّد الشافعي رحمه الله ببيان طلاقة القدرة الإلهية، وأن الله يفعل ما يشاء ويحكم ما يريد لا معقب لحكمه. ومن رحمته بعباده التدرج في التشريع؛ كتحريم الخمر الذي جاء متدرجاً، وكذلك الصيام. فنسخ بعض الأحكام وتغييرها هو من تمام رحمة الله بالعباد للتخفيف عنهم والتوسعة عليهم، كنسخ خمسين صلاة لتصبح خمساً في العمل وخمسين في الأجر. فالحمد لله على نعمه في إثبات الأحكام ونسخها.

\section{مذهب الشافعي في النسخ (السنة لا تنسخ القرآن)}
قرر الشافعي أصلاً عظيماً، وهو أن القرآن لا ينسخه إلا قرآن، والسنة لا ينسخها إلا سنة. ومذهبه أن السنة لا تنسخ القرآن أبداً، ولا يوجد مثال صحيح صريح سالم من المعارضة يثبت أن القرآن نُسخ بسنة.
وإنما السنة تابعة للكتاب، مفسرة ومبينة له. قال تعالى: ﴿قُلْ مَا يَكُونُ لِي أَنْ أُبَدِّلَهُ مِن تِلْقَاءِ نَفْسِي ۖ إِنْ أَتَّبِعُ إِلَّا مَا يُوحَىٰ إِلَيَّ﴾، فالنبي صلى الله عليه وسلم لا ينسخ حكم القرآن ولا يبدله من تلقاء نفسه، بل التشريع ونسخ القرآن من خصائص رب العالمين وحده.

وإذا نُسخت سنة بسنة، فإن النبي صلى الله عليه وسلم يوضح ذلك بسنة أخرى حتى تقوم الحجة، ولا يترك الأمة في لبس. وقد حذر الشافعي من دعوى أن السنة تُنسخ بالقرآن دون وجود سنة توضح ذلك، لئلا يتذرع أصحاب الأهواء (كالقرآنيين أو بعض المقلدين) برد السنن الصحيحة بحجة أنها تخالف القرآن أو أنها منسوخة، كما زعموا في أحكام البيوع المحرمة وحدود الزنا والمسح على الخفين.

\section{علاقة السنة بالقرآن والرد على المتعصبين}
السنة لا تكون أبداً إلا موافقة لكتاب الله عز وجل؛ لأن الكل وحي من الله. ومن ادعى التناقض بينهما فهو مشكك أو جاهل. واتباع السنة هو في الحقيقة اتباع للقرآن؛ لأن القرآن هو الذي أمر باتباعها، فمن قبل عن رسول الله فقد قبل عن الله أمره.

وقد نقل العلامة أحمد شاكر رحمه الله تحذيراً شديداً للمقلدين المتعصبين للمذاهب، الذين يردون الأحاديث الصحاح بدعوى احتمال نسخها نصرة لمذاهبهم (كقول بعضهم: كل حديث خالف مذهبنا فهو منسوخ أو مؤول). وهذا الصنيع يخرج عامة السنن من أيدي الناس، وهو الذي مهد الطريق لاحقاً للقوانين الوضعية والعلمانية لتحل محل الشريعة بدعوى المصلحة. ولا يسع مسلماً تبلغه سنة صحيحة إلا أن يتبعها.

\section{هل النسخ يكون إلى غير بدل؟}
قرر الشافعي، ووافقه جمهور المحققين كالعلامة الشنقيطي، أن النسخ لا يكون إلا إلى بدل، لقوله تعالى: ﴿نَأْتِ بِخَيْرٍ مِّنْهَا أَوْ مِثْلِهَا﴾. فليس ينسخ فرض أبداً إلا أثبت مكانه فرض، كنسخ التوجه لبيت المقدس وإثبات التوجه للكعبة مكانه. 

\section{أمثلة تطبيقية على الناسخ والمنسوخ}

\subsection{نسخ قيام الليل بالصلوات الخمس}
فرض الله قيام الليل أولاً بقوله: ﴿قُمِ اللَّيْلَ إِلَّا قَلِيلًا﴾، ثم نسخه وخففه بقوله: ﴿فَاقْرَءُوا مَا تَيَسَّرَ مِنْهُ﴾. ثم نُسخ ذلك كله وأصبح نافلة بقوله: ﴿وَمِنَ اللَّيْلِ فَتَهَجَّدْ بِهِ نَافِلَةً لَّكَ﴾. وبينت السنة المتواترة (كحديث الأعرابي) أن الله لم يفرض على العباد سوى الصلوات الخمس في اليوم والليلة، وأن ما عداها تطوع، ولا يجوز لأحد أن يقول بوجوب قيام الليل بعد هذا النسخ، وإن كان يستحب استحباباً شديداً.

\subsection{نسخ المصابرة في القتال}
قال تعالى: ﴿إِن يَكُن مِّنكُمْ عِشْرُونَ صَابِرُونَ يَغْلِبُوا مِائَتَيْنِ﴾، فكان الفرض ألا يفر الواحد من المسلمين أمام عشرة من المشركين. ثم خفف الله عنهم لعلمه بضعفهم، فنسخها بقوله: ﴿الْآنَ خَفَّفَ اللَّهُ عَنكُمْ وَعَلِمَ أَنَّ فِيكُمْ ضَعْفًا ۚ فَإِن يَكُن مِّنكُم مِّائَةٌ صَابِرَةٌ يَغْلِبُوا مِائَتَيْنِ﴾، فأصبح الفرض ألا يفر الواحد من المسلمين أمام اثنين من المشركين، والآية شارحة لنفسها ولا تحتاج لتفسير.

\subsection{نسخ حبس الزانيات}
كان حد الزنى في أول الأمر الحبس للمرأة والأذى للرجل: ﴿فَأَمْسِكُوهُنَّ فِي الْبُيُوتِ حَتَّىٰ يَتَوَفَّاهُنَّ الْمَوْتُ أَوْ يَجْعَلَ اللَّهُ لَهُنَّ سَبِيلًا﴾. فبين النبي صلى الله عليه وسلم هذا "السبيل" في حديث عبادة بن الصامت: «خذوا عني، قد جعل الله لهن سبيلا، البكر بالبكر جلد مائة وتغريب عام، والثيب بالثيب جلد مائة والرجم». ثم نُسخ الجلد عن الثيب بالرجم فقط كما فعل النبي صلى الله عليه وسلم مع ماعز والغامدية.

\section{أحكام فقهية مستنبطة (الحائض، المغمى عليه، السكران)}

\subsection{قضاء الصلاة للمرأة الحائض}
الحيض أمر جبلّي خلقه الله في المرأة، فهي ليست عاصية بوجوده. وقد أسقط الله عنها فرض الصلاة أيام حيضها. ولما كانت معذورة بغير جناية منها، أجمع العلماء على أنها لا تقضي الصلاة، ولكنها تقضي الصيام (لأن الصيام خفيف يقع شهراً في السنة، بخلاف الصلاة التي تتكرر يومياً فيشق قضاؤها).

\subsection{صلاة المغمى عليه والسكران}
المغمى عليه (أو المغلوب على عقله بمرض وعارض سماوي) لا يقضي الصلاة إذا خرج وقتها وهو فاقد للوعي؛ قياساً على الحائض، لأنه فقد مناط التكليف (العقل) بغير جناية منه ولا فعل.
أما "السكران"، فإذا فاتته الصلاة وجب عليه قضاؤها إذا أفاق، ولا يُقاس على المغمى عليه؛ لأن السكران أدخل المانع على نفسه متعمداً، فهو عاصٍ بفعله، والجهة غير منفكة، ولذا تبطل صلاته حال سكره لقوله تعالى: ﴿لَا تَقْرَبُوا الصَّلَاةَ وَأَنتُمْ سُكَارَىٰ﴾، ويلزمه القضاء بعد إفاقته عقوبة له.

\section{فقه إقامة الحدود}
ساق الشافعي حديث العسيف (الأجير) الذي زنى بامرأة مستأجره، فافتداه أبوه بمائة شاة وجارية ظناً منه أنه يُرجم، ثم سأل العلماء فاخبروه أن على ابنه جلد مائة وتغريب عام. فارتفعوا للنبي صلى الله عليه وسلم، فأبطل الفدية ورد الشياه والجارية، وأمر بجلد الشاب وتغريبه. ثم أرسل أُنَيْسًا الأسلمي للمرأة: «فإن اعترفت فارجمها».
وفي هذا دلالة عظيمة على كمال الشريعة؛ فإقامة الحدود لا تُبنى على الشبهات، بل النبي صلى الله عليه وسلم لم يكتفِ بإقرار الشاب والزوج، بل أرسل للمرأة لتُقر بنفسها، فلما اعترفت رُجمت، ولو أنكرت لَدُرئ عنها الحد. وهذا يرد على شبهات العلمانيين الذين يصورون الشريعة كأنها متعطشة للدماء، بل هي غاية في الرحمة والاحتياط للنفوس والأعراض.